% ============================================================
% METHODOLODY
% ============================================================
\section*{Methodology}

% ----- Excel Graph (2D) -----

\subsubsection*{Excel Graph (2D)}

Microsoft Excel is one of the simplest tools for generating graphs and calculating tabulated data. 
This approach is particularly useful for understanding preliminary concepts in aircraft design, especially wing profiles, 
which are fundamental components of an airplane. 

Using cosine spacing with 1000 points provides an efficient method for distributing coordinates along the chord. 
The primary advantage of cosine spacing is its ability to increase the density of points in regions with high curvature, 
thereby producing a smoother and more accurate airfoil representation.

There are two commonly used cosine spacing methods:

\begin{itemize}
    \item[(1)] The $x$-coordinates are defined using a sequence of angles ranging from $0$ to $\displaystyle \frac{\pi}{2}$. 
    This method concentrates more points near the leading edge, where the curvature is greatest, 
    while fewer points are distributed near the trailing edge.
    
    \item[(2)] The $x$-coordinates are defined using angles ranging from $0$ to $\pi$. This method distributes points symmetrically, 
    ensuring equal point density near both the leading edge and trailing edge. 
    However, the points near the mid-chord region are more sparsely distributed compared to the edge regions.
\end{itemize}

Since 1000 points are sufficient to produce a smooth graphical representation, 
both cosine spacing methods yield nearly identical results. 
In this study, the first method is adopted.

% ----- Airfoil Graphing -----

\subsubsection*{Airfoil Graphing}

For geometry-based thickness distributions, 
including the NACA 4-digit, 5-digit, and their modified series, 
the airfoil geometries are constructed using the analytical equations provided in the official NACA reports.
However, for airfoils designed based on thin airfoil theory, 
no explicit analytical thickness distributions are defined. 
Due to the complexity of aerodynamic analysis and the requirement for advanced computational tools, 
these airfoils are generated using available airfoil coordinate data. 
The governing equations are then approximated using Lagrange interpolation.

The geometric characteristics of airfoils vary significantly depending on changes in key parameters. 
These geometric variations directly influence aerodynamic performance, particularly lift, drag, and moment coefficients during flight.

In this study, 
one geometric parameter of a given NACA airfoil is varied while all other parameters are held constant. 
This approach allows for a clear observation of how individual parameters affect the airfoil shape.
